\chapter{Introducción}
\label{cap:introduccion}

\section{Planteamiento del problema}

La saturación hospitalaria y la asignación de recursos escasos son retos
persistentes del sistema de salud mexicano. Contar con herramientas que estimen
tempranamente el riesgo de un desenlace grave permite priorizar la atención y
anticipar necesidades de camas y ventiladores. La literatura reciente muestra,
sin embargo, que el problema puramente predictivo ya está ampliamente resuelto
para el caso mexicano: existen múltiples modelos con buen
desempeño~\parencite{becerra2022,mendez2024}.

El problema no abordado es otro. Un modelo puede tener excelente desempeño
promedio y, al mismo tiempo, ser sistemáticamente menos preciso para habitantes
de entidades con menor infraestructura, para personas mayores o para un sexo
determinado. Esta posibilidad no es hipotética: el caso mejor documentado de
sesgo algorítmico en salud es un modelo de gestión poblacional en Estados Unidos
que asignaba menos recursos a pacientes afroamericanos con igual carga de
enfermedad, no por usar la raza como predictor, sino porque el desenlace elegido
--el gasto sanitario-- estaba ya distorsionado por el acceso
desigual~\parencite{obermeyer2019}.

En un país tan heterogéneo como México, un modelo entrenado con datos dominados
por entidades urbanas puede desempeñarse peor justo donde el sistema de salud es
más débil, amplificando la inequidad en lugar de reducirla. El problema central
de esta tesis es, por tanto, diagnosticar y corregir la inequidad del modelo
--especialmente la geográfica-- sin sacrificar de forma inaceptable su utilidad
clínica.

\section{Preguntas de investigación}

\begin{enumerate}[label=P\arabic*., leftmargin=*]
  \item Dado un modelo de riesgo con buen desempeño global sobre datos
    mexicanos, ¿existen disparidades de desempeño entre entidades federativas,
    grupos etarios y sexos?
  \item ¿Cómo se relaciona la magnitud del sesgo geográfico con indicadores de
    infraestructura o marginación de cada entidad?
  \item ¿Qué técnicas de mitigación --pre, in y post-procesamiento-- reducen
    mejor esas disparidades, y cuál es el costo en desempeño global expresado
    como frontera de Pareto?
  \item ¿Se transfieren tanto el desempeño como la equidad del modelo a una
    fuente de datos independiente?
  \item De forma exploratoria, ¿aparecen inequidades adicionales en subgrupos
    interseccionales, por ejemplo adultos mayores de entidades marginadas?
\end{enumerate}

\section{Objetivos}

\subsection{Objetivo general}

Desarrollar, auditar y mitigar la equidad --con énfasis en el sesgo geográfico
entre entidades federativas-- de un modelo de estratificación de riesgo clínico
basado en datos abiertos de salud de México, valorando la transportabilidad de
sus hallazgos.

\subsection{Objetivos específicos}

\begin{enumerate}[label=OE\arabic*., leftmargin=*]
  \item Construir un canal reproducible de ingesta, limpieza y preparación de la
    base de la DGE, con control estricto de fuga de información.
  \item Entrenar un modelo de riesgo de referencia con desempeño comparable al
    estado del arte, como sujeto del análisis de equidad.
  \item Auditar la equidad del modelo por entidad federativa, sexo y grupo
    etario, y examinar la asociación del sesgo geográfico con indicadores
    contextuales.
  \item Aplicar y comparar técnicas de mitigación en las tres etapas del canal,
    construyendo la frontera de Pareto entre equidad y desempeño.
  \item Explorar la equidad en subgrupos interseccionales y traducir los
    hallazgos a recomendaciones de despliegue.
\end{enumerate}

\section{Hipótesis}

\begin{enumerate}[label=H\arabic*., ref=H\arabic*, leftmargin=*]
  \item \label{h1} Un modelo con buen desempeño global presentará disparidades
    de desempeño entre entidades federativas, medibles por \emph{equalized odds}
    y calibración por grupo.
  \item \label{h2} La magnitud del sesgo geográfico se asociará negativamente
    con la infraestructura de salud de cada entidad: peor desempeño donde hay
    menos recursos.
  \item \label{h3} Las técnicas de mitigación reducirán las disparidades con una
    pérdida de desempeño global inferior a un umbral predefinido, cuantificable
    mediante la frontera de Pareto.
  \item \label{h4} El desempeño se transferirá razonablemente a una fuente
    independiente, pero las disparidades podrían acentuarse, evidenciando la
    necesidad de recalibrado por contexto.
\end{enumerate}

\section{Contribuciones}

\begin{enumerate}[leftmargin=*]
  \item La primera auditoría de equidad geográfica --entidad federativa como
    atributo sensible-- de un modelo de riesgo clínico entrenado sobre datos
    abiertos mexicanos.
  \item Evidencia de que la disparidad geográfica no se explica por composición
    etaria: persiste al restringir el análisis al subgrupo de 60 años o más.
  \item Una comparación de las tres familias de mitigación bajo un punto de
    operación clínico común, que documenta una asimetría relevante: con 32
    grupos de tamaño muy desigual, solo el post-procesamiento por umbrales de
    grupo reduce la disparidad de manera apreciable.
  \item Una frontera de Pareto equidad--desempeño construida con umbrales
    estimados \emph{fuera de muestra}, y la cuantificación de cuánto optimismo
    introduce estimarlos en el propio conjunto de prueba.
  \item Evidencia de que la inequidad no es un problema de cobertura: un modelo
    que nunca vio una entidad se desempeña en ella casi igual que uno que sí la vio,
    y las brechas no se acentúan.
  \item Un canal reproducible con contrato de datos explícito entre etapas
    (apéndice~\ref{ap:reproducibilidad}).
\end{enumerate}

\section{Alcance y delimitaciones}
\label{sec:alcance}

Este documento reporta el trabajo correspondiente a las fases de datos, modelo
de referencia, auditoría de equidad y mitigación. Para que la lectura de los
resultados sea correcta, conviene fijar con precisión sus límites:

\begin{itemize}[leftmargin=*]
  \item \textbf{Cohorte.} Casos COVID confirmados
    (\texttt{CLASIFICACION\_FINAL} $\in \{1,2,3\}$) del corte anual 2021 de la
    base de la DGE, $n=\cohorteN$. No se emplea la serie histórica completa.
  \item \textbf{Desenlace.} Mortalidad. La hospitalización se deriva y se
    describe en el análisis exploratorio, pero no se modela en este documento.
    En esta cohorte fallecer implica estar registrado como hospitalizado, de modo
    que el desenlace solo es posible en el \fracHosp\,\% de los casos; la
    sección~\ref{sec:sensibilidad_res} analiza las consecuencias.
  \item \textbf{Validación externa.} La transportabilidad a la base de Egresos
    Hospitalarios de la DGIS, que sustenta \ref{h4}, \emph{no} se ha ejecutado, y
    \ref{h4} se reporta como trabajo pendiente (sección~\ref{sec:trabajo_futuro}).
    Sí se evalúa la generalización geográfica dentro de la misma fuente
    (sección~\ref{sec:loeo_res}), que no la sustituye.
  \item \textbf{Indicadores contextuales.} El examen de \ref{h2} se apoya en
    el índice de marginación por entidad de CONAPO 2020~\parencite{conapo2020}, que es
    un agregado socioeconómico estatal. No se emplearon indicadores de
    infraestructura de salud por unidad médica, que serían el \emph{proxy} más
    directo del mecanismo que \ref{h2} postula.
  \item \textbf{Interpretabilidad.} Se emplea importancia por permutación. No se
    calcularon valores SHAP ni análisis de curva de decisión.
  \item \textbf{Incertidumbre.} Las brechas del eje geográfico se acompañan de
    intervalos de confianza por remuestreo y de una prueba de permutación
    (sección~\ref{sec:incertidumbre}). Los subgrupos interseccionales, en cambio,
    se reportan como estimaciones puntuales.
\end{itemize}

\nota{Los límites anteriores no son concesiones retóricas. Cualquier afirmación
de esta tesis que dependa de \ref{h4} o de intervalos de confianza está marcada
como preliminar en el punto donde aparece.}
